% ============================================================
% SECCION 8 - METRICAS DE CALIDAD (Maria Escudero - PE5)
% Integra el fragmento M1-M4 de la auditoria estadistica
% (Robinson Jeanpierre) y desarrolla M5-M6 (Maria Escudero).
% Fragmento para \input en el informe final
% ============================================================
\section{Métricas de calidad}
\label{sec:metricas-calidad}

\subsection{Instrumento y conteos base}
\label{sec:conteos-base}

La auditoría aplica las seis métricas derivadas de ISO/IEC 25010:2023 definidas en la guía, cada una como fracción sobre un total contable, en dos olas de medición. La primera ola se midió sobre la línea base corregida del documento (42 requisitos: 27 funcionales RF-01 a RF-27 y 15 no funcionales RNF-01 a RNF-15), con conteos publicados antes de calcular en \texttt{09\_Cierre\_PE5/conteos\_base\_auditoria.csv}; entre ellos, 16 casos de uso identificados y especificados, 2 actores operativos, 16 historias de usuario y 861 pares de requisitos ($n(n-1)/2$ con $n=42$). Las métricas M1 a M4 fueron medidas por la auditoría estadística reproducible cuyo método, scripts y salidas constan en \texttt{09\_Cierre\_PE5/analisis\_estadistico/} (Subsección~\ref{subsec:m1-m4}); para M2 el análisis trabajó los 116 pares intra-ámbito conforme a las reglas documentadas en sus notas metodológicas. La segunda ola corresponde a la expansión del documento con los componentes de IA (RF-28 a RF-31 y RNF-16 a RNF-26, hasta 57 requisitos); los conteos actualizados están en \texttt{conteos\_base\_auditoria\_actualizado.csv}. La re-auditoría sobre esa versión detectó en su primera medición 54/57 atributos completos ($94{,}7~\%$), por Origen sin fuente trazable en tres fichas de equidad; aplicada la corrección sobre RNF-20, RNF-24 y RNF-25 en \texttt{ERS\_SRS\_2A\_v2.0.tex} y re-medido, dio 57/57. En consistencia, la segunda ola midió 191 pares intra-ámbito en nueve ámbitos sin conflictos; el total teórico todos-contra-todos para $n=57$ es 1596 y queda registrado solo como conteo base. M5 y M6 se midieron sobre la línea base corregida de 42 requisitos en la fecha de la re-inspección, con la muestra y el denominador declarados explícitamente en cada caso.

\subsection{Tabla de auditoría}
\label{sec:tabla-auditoria}

\begin{center}
\begin{longtable}{|p{2.4cm}|p{3.0cm}|p{2.2cm}|p{1.9cm}|p{1.2cm}|p{2.8cm}|}
\hline
\textbf{Métrica} & \textbf{Valor obtenido} & \textbf{Referencia} & \textbf{Norma} & \textbf{Cumple} & \textbf{Acción de mejora} \\
\hline
\endfirsthead
\hline
\textbf{Métrica} & \textbf{Valor obtenido} & \textbf{Referencia} & \textbf{Norma} & \textbf{Cumple} & \textbf{Acción de mejora} \\
\hline
\endhead
\hline
\endfoot
\hline
\endlastfoot
M1 Completitud & M1a = 100~\% (ola 2: 57/57 tras corrección desde $94{,}7~\%$); M1b = N/D en este repositorio; M1c = 100~\% (2/2 operativos) & $\geq$95~\%; $=$100~\%; $=$100~\% & ISO 25010; 29148 & Sí (M1b con causa y acción) & Corrección de origen en tres fichas y verificación cruzada de M1b contra el catálogo de casos de uso del equipo \\
\hline
M2 Consistencia & $1 - 0/191 = 1{,}00$ (pares intra-ámbito, nueve ámbitos; ola 1: $1 - 0/116$) & $\geq$0,98 y cero conflictos abiertos & ISO 25010 & Sí & Observaciones de solapamiento registradas sin conflicto en ambas olas \\
\hline
M3 Verificabilidad & 100~\% (57/57 con umbral, unidad y método; ola 1: 42/42) & $\geq$90~\% & ISO 29148 & Sí & Recomendación menor de redacción en RNF-07, RNF-09 y RNF-11 \\
\hline
M4 Trazabilidad & Adelante 16/16 en alcance Debe tener = 100~\%; atrás 57/57 = 100~\% (ola 1: 42/42) & $\geq$90~\%; $=$100~\% & ISO 29148 & Sí & Tramo CU$\rightarrow$clase$\rightarrow$CP pendiente de verificación cruzada \\
\hline
M5 Modificabilidad & 3,50 antes; 3,00 después & $\leq$3,0 & ISO 25010 & Sí tras corrección & Reespecificación de vistas (Subsección~\ref{sec:m5-modificabilidad}) \\
\hline
M6 Corrección & 0,119 antes; 0,024 después & $\leq$0,05 & ISO 25010 & Sí tras corrección & Correcciones DR-02 a DR-04 (Sección~\ref{sec:correcciones-reinspeccion}) \\
\hline
\end{longtable}
\label{tab:auditoria-metricas}
\end{center}

\noindent La única componente sin valor al cierre es M1b, porque el catálogo completo de casos de uso vive en el modelado del equipo y no en este repositorio; se declara como no medible con causa y acción registradas, conforme al formato del instrumento. Para M1c se cuentan como actores los dos actores operativos del sistema (Administrador y Trabajador agrícola), ambos con al menos un requisito funcional asociado; las otras dos partes interesadas del mapa (Equipo de desarrollo y Docente) no interactúan con el sistema, lectura justificada con el modelo i* y documentada en \texttt{m1\_completitud.md}.

% Fragmento para la seccion 8 (Metricas de calidad) del informe final.
% Contiene las metricas M1 a M4 medidas por la auditoria estadistica en dos olas.
% Integrar junto a sec8_metricas.tex, que cubre M5 y M6.

\subsection{Auditoría estadística de completitud, consistencia, verificabilidad y trazabilidad}
\label{subsec:m1-m4}

Las métricas M1 a M4 fueron medidas por conteo directo mediante el script
reproducible \texttt{conteo\_requisitos.ps1}, cuyas salidas CSV están
versionadas en el repositorio. Se midió en dos olas: la primera sobre la línea
base corregida (\texttt{ERS\_SRS\_2A\_v1.0.tex}, 42 requisitos) y la segunda
sobre la versión fusionada con los componentes de IA
(\texttt{ERS\_SRS\_2A\_v2.0.tex}, 57 requisitos), seleccionada con el
parámetro \texttt{-ErsPath}. Los conteos base se publicaron antes de calcular
cada métrica y una segunda ejecución produjo salidas idénticas (verificación
por hash SHA-256). La tabla reporta los valores finales de la segunda ola; los
valores de la primera se citan entre paréntesis cuando difieren.

\begin{center}
\begin{longtable}{|p{3.2cm}|p{2.8cm}|p{2.2cm}|p{1.5cm}|p{4.0cm}|}
\hline
\textbf{Métrica} & \textbf{Valor obtenido} & \textbf{Referencia} & \textbf{Cumple} & \textbf{Acción} \\
\hline
\endfirsthead
\hline
\endhead
M1a Completitud de atributos & $54/57 \times 100 = 94{,}7~\% \rightarrow 57/57 = 100~\%$ (ola 1: $42/42$) & $\geq 95~\%$ & Sí (tras corrección) & Primera medición detectó Origen sin fuente trazable en RNF-20/24/25; añadida la referencia al requisito de origen y re-medido \\
\hline
M1b CU especificados & No medible en este repositorio & $= 100~\%$ & N/D & Verificación cruzada contra el catálogo de CU del equipo \\
\hline
M1c Actores operativos con $\geq 1$ RF & $2/2 \times 100 = 100~\%$ & $= 100~\%$ & Sí & Lectura alternativa 2/4 justificada vía i* \\
\hline
M2 Consistencia & $1 - 0/191 = 1{,}00$ (ola 1: $1 - 0/116$) & $\geq 0{,}98$ & Sí & Nueve ámbitos en la ola 2; observaciones de solapamiento registradas, sin conflicto \\
\hline
M3 Verificabilidad & $57/57 \times 100 = 100~\%$ (ola 1: $42/42$) & $\geq 90~\%$ & Sí & Las 15 fichas de IA nacieron con umbral, unidad y método \\
\hline
M4ade Trazabilidad adelante & $16/16 \times 100 = 100~\%$ (alcance Debe tener) & $\geq 90~\%$ & Sí & RF-28 a RF-31 son Debería tener; tramo CU$\rightarrow$clase$\rightarrow$CP pendiente de verificación cruzada \\
\hline
M4atr Trazabilidad atrás & $57/57 \times 100 = 100~\%$ (ola 1: $42/42$) & $= 100~\%$ & Sí & Cerrado junto con M1a (mismo atributo Origen) \\
\hline
\end{longtable}
\label{tab:auditoria-m1-m4}
\end{center}

El análisis de consistencia examinó 116 pares intra-ámbito en la primera ola
(42 requisitos, ocho ámbitos funcionales) y 191 pares en la segunda (57
requisitos, nueve ámbitos), sin conflictos directos ni indirectos abiertos en
ninguna de las dos. La auditoría demostró además el ciclo completo de mejora:
la incorporación de los 15 requisitos de IA había degradado M1a a $94{,}7~\%$
por tres fichas sin fuente trazable; localizada la causa, la corrección se
aplicó sobre las fichas de equidad (RNF-20, RNF-24 y RNF-25) y la re-medición
confirmó el $100~\%$. La prueba de verificabilidad aplicó las cuatro preguntas
de la guía a cada requisito; el escaneo de vocabulario no medible arrojó cero
coincidencias. En trazabilidad, los 16 requisitos de prioridad «Debe tener»
cuentan con historia de usuario y escenario de aceptación Gherkin, lo que
arroja una sincronización backlog--ERS del $100~\%$; no se detectaron
requisitos huérfanos ni historias sin requisito de origen.

Los conteos completos, la matriz de pares, la aritmética por métrica con ambas
olas y la cadena por requisito constan en el directorio
\texttt{09\_Cierre\_PE5/analisis\_estadistico/} del repositorio.


\subsection{M5 --- Modificabilidad}
\label{sec:m5-modificabilidad}

Regla operativa declarada: un requisito se cuenta como afectado al modificar otro cuando su ficha lo referencia explícitamente en entradas, salidas, precondiciones o criterio de verificación. La tabla completa de dependencias, con la línea del documento que respalda cada enlace, está en \texttt{09\_Cierre\_PE5/calculo\_M5\_modificabilidad.csv}. La muestra de seis requisitos cubre los módulos de acceso, tareas, cosechas, inventario, registro fitosanitario e IA:

\begin{center}
\begin{longtable}{|p{2.0cm}|p{9.0cm}|p{1.7cm}|}
\hline
\textbf{Modificado} & \textbf{Requisitos afectados} & \textbf{Cantidad} \\
\hline
\endfirsthead
\hline
\textbf{Modificado} & \textbf{Requisitos afectados} & \textbf{Cantidad} \\
\hline
\endhead
\hline
\endfoot
\hline
\endlastfoot
RF-01 & RF-14 (permisos por rol); RNF-10 (bloqueo sobre inicio de sesión) & 2 \\
\hline
RF-04 & RF-10 (consulta de actividades); RF-24 (historial incluye labores) & 2 \\
\hline
RF-05 & RF-03, RF-07, RF-11, RF-12, RF-17, RF-18, RF-24, RF-27 & 8 \\
\hline
RF-06 & RF-21 (insumo del calendario); RF-22 (descuento de inventario); RF-25 (alerta por disponibilidad crítica) & 3 \\
\hline
RF-16 & RF-03, RF-19, RF-20, RF-25 & 4 \\
\hline
RF-26 & RF-15 (mismo módulo de IA); RNF-15 (explicabilidad de la alerta temprana) & 2 \\
\hline
\end{longtable}
\label{tab:m5-muestra}
\end{center}

\noindent Aritmética de la medición inicial:
\[ M5 = \frac{2+2+8+3+4+2}{6} = \frac{21}{6} = 3{,}50 \]
El valor excede la referencia ($\leq$ 3,0). El recorrido muestra que el acoplamiento se concentra en RF-05, punto único de captura de producción consumido directamente por cuatro consultas y reportes.

Acción de mejora aplicada: RF-12, RF-18 y RF-27 se reespecifican como vistas derivadas de la consulta de producción por período (RF-07), de modo que consumen su resultado en lugar de los registros crudos de RF-05. Desaparecen tres enlaces directos hacia RF-05 y la re-medición sobre la misma muestra da:
\[ M5' = \frac{2+2+5+3+4+2}{6} = \frac{18}{6} = 3{,}00 \]
El valor cumple la referencia. Queda registrada como recomendación no aplicada la fusión eventual de RF-17 dentro de RF-05 como atributos opcionales del registro de cosecha, que reduciría el acoplamiento medio a $17/6 \approx 2{,}83$.

\subsection{M6 --- Corrección}
\label{sec:m6-correccion}

El numerador corresponde a los defectos residuales hallados en la re-inspección del 22-08-2026 sobre el ERS ya corregido entre entregas (Sección~\ref{sec:validacion}); el denominador es el total de requisitos de esa línea base (42). Medición inicial, sobre los cinco hallazgos del registro:
\[ M6_{antes} = \frac{5}{42} = 0{,}119 \]
El valor excede la referencia ($\leq$ 0,05). Tras aplicar las correcciones DR-02 a DR-04 y cerrar documentalmente DR-05, la re-medición sobre defectos abiertos da:
\[ M6_{despues} = \frac{1}{42} = 0{,}024 \]
El valor cumple la referencia. El residual restante (DR-01, filas de la matriz) quedó cerrado el 23-08-2026 dentro del Paso 2 con la ampliación de la matriz de trazabilidad a 62 filas de datos con las diez columnas de la plantilla, con lo que el residual final del ciclo es $0/42 = 0{,}00$; el par oficial de la corrección se mantiene en 0,119 $\rightarrow$ 0,024 por ser el medido al cierre del ciclo documental.

\subsection{Comparación global}
\label{sec:comparacion-global}

Al cierre de la práctica las seis métricas alcanzan su valor de referencia: cuatro cumplieron en la primera medición y dos requirieron corrección documental aplicada durante la práctica (M5 y M6), ambas re-medidas con aritmética visible y par antes/después registrado. La única componente declarada sin valor es M1b, por perímetro del repositorio y con causa y acción registradas. La auditoría de la segunda ola ejercitó además el ciclo completo de mejora sobre la expansión a 57 requisitos: detectó una caída puntual de completitud ($94{,}7~\%$), la atribuyó a tres fichas concretas, aplicó la corrección y la re-medición confirmó el $100~\%$, de modo que el cierre quedó sin degradación de los indicadores.
